\documentclass[11pt,a4paper]{article}
\usepackage[margin=1in]{geometry}
\usepackage{multicol}
\usepackage{amsmath,amsfonts,amssymb}
\usepackage{array}
\usepackage[caption=false,font=normalsize,labelfont=sf,textfont=sf]{subfig}
\usepackage{textcomp}
\usepackage{url}
\usepackage[hidelinks]{hyperref}
\usepackage{verbatim}
\usepackage{graphicx}

\begin{document}

\title{The Science Genome: A Framework for Semantic Scientific Analysis}

\author{Louis~George}

\maketitle

\begin{abstract}
This paper presents a standalone framework for analyzing how scientific ideas are inherited, transformed, and recombined across publication lineages. Publications are represented in a shared semantic space and decomposed into inherited content explained by cited predecessors plus residual contribution not explained by those predecessors. The framework combines constrained inheritance inference on a citation graph, enabling interpretable edge-level and node-level analysis of knowledge flow. The result is specification for content-aware scientometric analysis that supports lineage tracing, novelty assessment, and reproducible research interpretation.
\end{abstract}

\begin{multicols}{2}

\section{Introduction}
Scientific progress emerges through reuse, adaptation, and recombination of prior ideas. Citation networks record part of this process, but citation counts and centrality metrics alone do not reveal how intellectual content flows, or how contributions from  predecessors impact those in newer papers.

This work formalises a content-aware framework for that distinction. Each paper is represented in a shared semantic space and approximated as a constrained mixture of cited-parent embeddings plus a residual term. The mixture weights encode lineage-specific inheritance, while the residual captures contribution not attributable to cited predecessors.

\subsection{Problem focus}
The framework is designed to support analysis of:
\begin{itemize}
    \item estimating parent-specific inheritance strength for each citation edge,
    \item quantifying paper-level residual contribution not explained by cited parents,
    \item mapping branch-level continuity and separation from aggregated edge and residual signals,
    \item and measuring cross-field transfer through inheritance mass that crosses discipline labels.
\end{itemize}

\subsection{Research stance}
The objective is analytical interpretation, not causal proof of influence. Inheritance and contribution factors are evidential signals that require uncertainty-aware reading and expert contextualization.

\subsection{Contributions of this paper}
This paper delivers a final, standalone specification by providing:
\begin{itemize}
    \item a formal decomposition model of parent-mixture inheritance and residual contribution,
    \item a constrained estimation procedure with uncertainty reporting for stable edge-level attribution,
    \item and a verification--validation protocol that tests specification compliance and interpretive utility.
\end{itemize}


\section{Related Work}\label{sec:related}

[Is there any semantic analysis of science papers?]
There are two types of analysis that can be done on scientific publications, semantic and relational.

\subsection{Citation Analysis} % Relational Analysis?
This is a study that finds patterns in different citation links to get insights such as mapping influence.\\*
These come in the form of 1) Link-Structure Rankings, 2) Pattern-Based Similarity, 3) Community \& Role Structure, 4) Temporal \& Evolutionary Citation Analytics, 5) Field-Normalisation \& Comparability, 6) Author/Entity-Level Indices, and 7) Counting \& Data Hygiene. A brief dive into each of them before a reflection from a GFSP perspective is given.

\subsubsection{Link-Structure Rankings}
\subsubsection{Pattern-Based Similarity}
\subsubsection{Community \& Role Structure}
\subsubsection{Temporal \& Evolutionary Citation Analytics}
\subsubsection{Field-Normalisation \& Comparability}
\subsubsection{Author/Entity-Level Indices}
\subsubsection{Counting \& Data Hygiene }


\subsubsection{Reflection}
These tend to miss semantic context, which is important as this holds the actual concepts, ideas, methods, and knowledge. They also hold no meaning or weightings in individual citations which induces a lot of noise that is only removed with large samples samples.


\subsection{Semantic Analysis}
[Intro words to semantic analysis]
This will explore: 1) Representation Foundations, 2) Semantic Similarity \& Mapping, 3) Concept \& Entity Normalisation, 4) Citation-Context Semantics, 5) Temporal \& Evolutionary Semantics, 6) Text-Graph Hybrids, and 7) Evaluation, Data, \& Hygiene.

\subsubsection{Representation Foundations}
\subsubsection{Semantic Similarity \& Mapping}
\subsubsection{Concept \& Entity Normalisation}
\subsubsection{Citation-Context Semantics}
\subsubsection{Temporal \& Evolutionary Semantics}
\subsubsection{Text-Graph Hybrids}
\subsubsection{Evaluation, Data, \& Hygiene}

\subsection{Methodological \& Theoretical Analysis}


\subsection{Traditional Metrics}
Citations, h-index, impact factor.

\textbf{Pros}: simple, widely used.

\textbf{Cons}: popularity bias, prestige effects, salami-slicing.

\subsection{Network-based Metrics}
PageRank, influence measures.

Better than counts, but still blind to semantic novelty.

Altmetrics \& Social Indicators: Twitter, Mendeley, etc.

Capture visibility, but amplify hype rather than intellectual content.

\subsection{Semantic Scientometrics}
Embeddings (topic models, SciBERT, SBERT, cosine similarity).

Start of capturing meaning, but usually measure similarity not inheritance/novelty.

\subsection{Paradigm Shift Detection}
Kuhn’s theory, attempts to quantify scientific revolutions.

Good framing, but not operationalised at semantic level.

\subsection{Gap Statement}
None of the above frameworks separates inherited vs novel contribution → motivation for your genetic evolutionary model.


\section{Method}\label{sec:method}
The method has four stages: (1) construct a semantically embedded paper corpus, (2) build a citation-based DAG, (3) estimate inheritance weights from cited parents to each focal paper, and (4) derive contribution and lineage statistics from the fitted decomposition.

\subsection{Construction of science space}
Let each paper $p_i$ be represented by title and abstract text. We encode this text using an appropriate transformer (e.g.  SciBERT) to obtain an embedding $\mathbf{e}_i \in \mathbb{R}^d$. Embeddings are $\ell_2$-normalized to stabilize cosine geometry and reduce scale effects between documents of different length (large documents having large magnitudes and seen as "more simular").

To improve robustness, we apply light preprocessing (Unicode normalization, citation-marker removal, and deduplication by identifier).

\subsection{Construction of citation DAG}
From citation metadata, we create a directed edge $(j \rightarrow i)$ when paper $p_i$ cites paper $p_j$. For each focal paper, we define its parent set $\mathcal{P}(i)$ as cited predecessors retained after filtering. We then enforce acyclicity by removing chronology-violating edges where publication year$(p_j)$ exceeds publication year$(p_i)$.

\subsection{Inheritance inference}
For each eligible paper $p_i$, we model its embedding as
\begin{equation}
\mathbf{e}_i \approx \sum_{j \in \mathcal{P}(i)} w_{ij}\,\mathbf{e}_j + \mathbf{r}_i,
\end{equation}
where $0 \le w_{ij} \le 1$ and $\mathbf{r}_i$ is a residual contribution vector.

Weights are estimated with non-negative constrained least squares and simplex regularization:
\begin{equation}
\min_{\mathbf{w}_i \ge 0}\; \left\|\mathbf{e}_i-\mathbf{E}_{\mathcal{P}(i)}\mathbf{w}_i\right\|_2^2 + \lambda\|\mathbf{w}_i\|_2^2
\quad\text{s.t.}\quad
\sum_{j \in \mathcal{P}(i)} w_{ij} \le 1,
\end{equation}
with $\mathbf{E}_{\mathcal{P}(i)}$ the parent-embedding matrix.

The non-negativity constraint enforces interpretable additive inheritance, and limiting the total parent weight to at most 1 prevents over-attribution when parent embeddings are highly collinear. The unassigned mass,
\begin{equation}
c_i = 1 - \sum_{j \in \mathcal{P}(i)} w_{ij},
\end{equation}
serves as a scalar \emph{contribution factor} (higher implies less explainability from cited parents). The vector residual
\begin{equation}
\mathbf{r}_i = \mathbf{e}_i-\sum_{j \in \mathcal{P}(i)} w_{ij}\mathbf{e}_j
\end{equation}
captures direction-specific novelty for downstream semantic inspection.

\subsection{Handling non-orthogonal parents}
Cited parents are typically correlated, which can make individual inheritance weights unstable even when overall reconstruction is good. We therefore use ridge regularization ($\lambda>0$) to reduce variance in coefficient allocation and improve numerical stability. We additionally apply bootstrap re-estimation over parent subsets to quantify uncertainty in $w_{ij}$, reporting confidence intervals alongside point estimates to distinguish robust inheritance edges from ill-conditioned assignments.

\subsection{Derived metrics and analyses}
From fitted weights we compute:
\begin{itemize}
    \item \textbf{Edge inheritance strength} $w_{ij}$ for lineage tracing.
    \item \textbf{Paper contribution and contribution  magnitude} $\mathbf{r}_i$, and  $c_i$ or  $\|\mathbf{r}_i\|_2$. % Have some thinking time of what c_i and ||r_i||_2  individually represent. 
    \item \textbf{Lineage concentration} (entropy of $\mathbf{w}_i$), indicating whether inheritance is diffuse or dominated by a few parents.
    \item \textbf{Cross-field transfer} by aggregating inheritance mass across discipline labels.
\end{itemize}

These quantities support both ranking-style comparisons and structural analyses of knowledge evolution in the citation DAG.


\section{Evaluation Protocol}\label{sec:experiments}
\subsection{Experimental objective}
The experiments are designed to test the framework claims directly: that citation-conditioned semantic decomposition yields interpretable inheritance edges, meaningful residual contribution signals, and robust conclusions under perturbation.

\subsection{Claim-driven evaluation skeleton}
\begin{enumerate}
    \item \textbf{Claim A: semantic reconstruction is informative rather than trivial.}
    \begin{itemize}
        \item \emph{Evidence target}: child embeddings are reconstructed by cited-parent mixtures with materially better fidelity than naive baselines.
        \item \emph{Suggested tests}: reconstruction error against (i) uniform parent weights, (ii) random cited-parent subsets, and (iii) popularity-weighted mixtures.
        \item \emph{Suggested metrics}: mean squared reconstruction error, cosine reconstruction score, and rank stability of high-contribution papers under metric choice.
    \end{itemize}

    \item \textbf{Claim B: inferred inheritance edges are numerically stable and interpretable.}
    \begin{itemize}
        \item \emph{Evidence target}: high-weight edges persist under parent-set perturbation and do not collapse under mild regularization changes.
        \item \emph{Suggested tests}: bootstrap re-estimation over parent subsets; sensitivity sweep over $\lambda$.
        \item \emph{Suggested metrics}: edge-weight confidence intervals, top-$k$ edge overlap across runs, and variance of node-level lineage concentration.
    \end{itemize}

    \item \textbf{Claim C: residual contribution captures signal distinct from citation popularity.}
    \begin{itemize}
        \item \emph{Evidence target}: contribution ranking is not reducible to citation count, venue prestige, or recency.
        \item \emph{Suggested tests}: correlation and partial-correlation analyses against popularity-oriented indicators; matched-group comparisons.
        \item \emph{Suggested metrics}: Spearman/Pearson correlation, explained variance gap versus baselines, and divergence profiles of top-ranked papers.
    \end{itemize}

    \item \textbf{Claim D: representation quality supports downstream inference.}
    \begin{itemize}
        \item \emph{Evidence target}: science-space structure is coherent enough for inheritance decomposition to be meaningful.
        \item \emph{Suggested tests}: nearest-neighbor density summaries and domain-tagged low-dimensional maps for semantic clustering.
        \item \emph{Suggested metrics}: density-distribution summaries, within-domain versus cross-domain neighborhood similarity, and cluster-separation diagnostics.
    \end{itemize}
\end{enumerate}

\subsection{Minimum evidence package for the draft}
\begin{itemize}
    \item Reconstruction-fidelity table with baseline comparisons.
    \item Inheritance-stability table (bootstrap confidence intervals and overlap statistics).
    \item Contribution-versus-popularity divergence analysis.
    \item Domain-tagged semantic map and nearest-neighbor density summary.
\end{itemize}

\subsection{Failure and limitation disclosure template}
\begin{itemize}
    \item Missing or noisy citation edges that alter parent sets.
    \item Embedding degradation for domain-specific language and terminology drift.
    \item Collinearity-driven instability in weak inheritance edges.
    \item Interpretation risk where residual signal may conflate novelty and representation artifacts.
\end{itemize}


\subsection{Implementation mapping for reproducible evidence}
The claim-linked evidence package is implemented as callable code paths:
\begin{itemize}
    \item Embedding diagnostics and density summaries: \texttt{src.analysis.structure.embedding\_norm\_diagnostics} and \texttt{src.analysis.structure.nearest\_neighbor\_density}.
    \item Solver convergence and conditioning summaries: \texttt{src.analysis.validation.generate\_validation\_report} (solver section).
    \item Uncertainty outputs (confidence-interval width and selection frequency): \texttt{src.analysis.validation.generate\_validation\_report} (uncertainty section).
    \item Lineage sanity checks: \texttt{src.analysis.validation.generate\_validation\_report} (lineage\_sanity section).
\end{itemize}
This mapping ensures each narrative claim can be traced to a deterministic programmatic report artifact.


\section{Discussion}\label{sec:discussion}
The pseudo-genetic formulation reframes evaluation from visibility-only indicators toward compositional analysis of knowledge flow. By separating inherited semantic mass from residual contribution, the framework provides interpretable evidence at both citation-edge and paper levels.

\subsection{How to read contribution and inheritance}
A high contribution factor indicates lower explainability by cited parents under the selected representation; it is not a direct claim of quality or impact. A high inheritance factor indicates stronger semantic dependence on predecessors; it is not a deficit signal. Synthesis papers, benchmarks, and consolidation work can have high scientific value despite high inherited mass.

\subsection{Analytical uses}
The framework supports:
\begin{itemize}
    \item lineage tracing for hypothesis-driven analysis,
    \item comparison of inheritance concentration across branches,
    \item detection of potential cross-domain transfer channels,
    \item and structured review of novelty claims against citation-supported dependence.
\end{itemize}

\subsection{Interpretive boundaries}
Outputs should not be used as deterministic judgments of researcher merit and should not be treated as causal evidence of intellectual influence without corroborating analysis. Strong claims should rely on convergence between quantitative outputs, uncertainty summaries, and expert qualitative inspection.


\section{Limitations}\label{sec:lim_future}
\subsection{Core limitations}
The framework has structural limits that must be accounted for in interpretation:
\begin{itemize}
    \item Embeddings compress complex text and may miss methodological detail.
    \item Citation graphs are incomplete proxies for true intellectual dependence.
    \item Correlated parent sets create attribution non-uniqueness without regularization.
    \item Residual contribution may reflect novelty, citation omission, or representation error.
\end{itemize}

\subsection{Reporting requirements}
Responsible use requires explicit reporting of:
\begin{itemize}
    \item confidence or stability summaries for inheritance estimates,
    \item preprocessing and parameterization choices,
    \item distinction between descriptive evidence and causal interpretation,
    \item and sensitivity of conclusions to weak or unstable edges.
\end{itemize}

These requirements are part of the framework definition and are necessary for publishable, reproducible analysis.


\section{Conclusion}\label{sec:conclusion}
This paper presented a pseudo-genetic framework for analyzing scientific publications in which each paper is decomposed into inherited semantic content from cited predecessors and a residual contribution component. By combining SciBERT embeddings, constrained inheritance inference, and citation-DAG analysis, the approach provides interpretable edge-level and node-level signals of knowledge flow.

Empirical results on a medium-scale open corpus show that inheritance and contribution factors capture information not contained in citation counts or network centrality alone. The method highlights selective parent dependence, surfaces plausible high-novelty papers that are not necessarily highly cited, and reveals cross-disciplinary transfer pathways.

Overall, the framework advances scientometric analysis toward content-aware evaluation of contribution. It is designed as a foundation for broader studies of scientific evolution, including larger corpora, richer full-text representations, and prospective monitoring of emerging research fronts.


\end{multicols}

\bibliographystyle{plain}
\bibliography{references}

\section*{TEMP: Todo}

\begin{itemize}
    \item Go through current paper, making sure everything I currently want in it is in it, and that it is written from the perspective of the paper, not from the perspective of my task (I am sus about the V\&V parts being in there)
    \item Make implementation lean and suplimentary to paper. Run investigations for the testing and validation
    \item Update the paper with what was seen in the testing and validation, any additional maths, etc etc
    \item Add diagrams to help explian concepts and visualise results (the topic map, inherritence etc) - make sure to only add things that add, keep them minimal and professional looking, and remember the intent - it isn't just to make it look fancy, it is to supliment the reading and make the concepts more intuitive.
    \item Add references where needed - if I make a claim that isn't mine, or reference a prior idea, add a reference to it
    \item Review for next steps
\end{itemize}


\end{document}
